\documentclass[twocolumn,10pt,letterpaper]{article}

% -- Page geometry --
\usepackage[
  top=0.85in,
  bottom=0.75in,
  left=0.65in,
  right=0.65in,
  columnsep=0.25in
]{geometry}

% -- Beautiful typography: Century Schoolbook --
\usepackage[T1]{fontenc}
\usepackage{tgschola}              % TeX Gyre Schola (Century Schoolbook)
\usepackage[scaled=0.82]{beramono} % Bera Mono for code
\usepackage{microtype}             % Microtypographic refinements
\usepackage{setspace}              % Line spacing control
\setstretch{1.08}                  % Slightly open leading for readability

% -- Packages --
\usepackage{xurl}        % Allow URL breaks at any character (load before hyperref)
\usepackage{hyperref}
\usepackage{graphicx}
\usepackage{booktabs}    % Professional tables
\usepackage{tabularx}    % Width-constrained tables
\usepackage{enumitem}    % List customization
\usepackage{fancyhdr}
\usepackage{xcolor}
\usepackage{balance}
\usepackage{stfloats}    % Enable [h] placement for table* in two-column layout
\usepackage{titlesec}    % Section heading customization
\usepackage{listings}    % Code blocks with line wrapping

% -- Paragraph spacing (no indent, modest skip) --
\setlength{\parindent}{0pt}
\setlength{\parskip}{0.4em plus 0.1em minus 0.05em}
\setlength{\emergencystretch}{1em}  % Extra stretch to avoid overfull hboxes in narrow columns

% -- Code listings with line wrapping --
\lstset{
  basicstyle=\small\ttfamily,
  breaklines=true,
  breakatwhitespace=false,
  breakautoindent=true,
  postbreak=\mbox{\textcolor{accentrule}{$\hookrightarrow$}\space},
  columns=fullflexible,
  keepspaces=true,
  frame=none,
  aboveskip=0.5em,
  belowskip=0.5em,
  xleftmargin=0pt,
  xrightmargin=0pt,
}

% -- Colors --
\definecolor{linkblue}{HTML}{1D4ED8}
\definecolor{headergray}{HTML}{1F2937}
\definecolor{rulegray}{HTML}{D1D5DB}
\definecolor{titlebg}{HTML}{111827}
\definecolor{accentrule}{HTML}{6B7280}
\definecolor{brandred}{HTML}{8B2635}

% -- Hyperlinks --
\hypersetup{
    colorlinks=true,
    linkcolor=linkblue,
    urlcolor=linkblue,
    citecolor=linkblue,
    pdfstartview=FitH
}

% -- Section numbering: Roman > Alph > Arabic > alph, no decimals --
\renewcommand{\thesection}{\Roman{section}}
\renewcommand{\thesubsection}{\Alph{subsection}}
\renewcommand{\thesubsubsection}{\arabic{subsubsection}}
\renewcommand{\theparagraph}{\alph{paragraph}}
\makeatletter
\@addtoreset{subsection}{section}
\@addtoreset{subsubsection}{subsection}
\@addtoreset{paragraph}{subsubsection}
\makeatother

% -- Section headings --
\titleformat{\section}
  {\large\bfseries\color{titlebg}}
  {\thesection.}{0.5em}{}
  [\vspace{-0.3em}{\color{accentrule}\rule{\columnwidth}{0.5pt}}]
\titleformat{\subsection}
  {\normalsize\bfseries\color{headergray}}
  {\thesubsection.}{0.4em}{}
\titleformat{\subsubsection}
  {\small\bfseries\itshape\color{headergray}}
  {\thesubsubsection.}{0.4em}{}
\titleformat{\paragraph}[runin]
  {\small\bfseries\color{headergray}}
  {\theparagraph)}{0.4em}{}[.\hspace{0.4em}]
\titlespacing*{\section}{0pt}{1.0em}{0.35em}
\titlespacing*{\subsection}{0pt}{0.75em}{0.2em}
\titlespacing*{\subsubsection}{0pt}{0.55em}{0.15em}
\titlespacing*{\paragraph}{0pt}{0.4em}{0em}

% -- Lists: compact, elegant --
\setlist{nosep, leftmargin=1.1em, topsep=0.2em, itemsep=0.05em}
\setlist[itemize]{label={\small\textcolor{accentrule}{\textbullet}}}
\setlist[enumerate]{label={\small\textcolor{headergray}{\arabic*.}}}

% -- Tables: tighter, small text --
\renewcommand{\arraystretch}{1.15}

% -- Header/footer (pages 2+) --
\pagestyle{fancy}
\fancyhf{}
\renewcommand{\headrulewidth}{0.3pt}
\renewcommand{\headrule}{\hbox to\headwidth{\color{rulegray}\leaders\hrule height \headrulewidth\hfill}}
\fancyhead[L]{\footnotesize\color{headergray}\textit{Codebase Review: beads\textbackslash{}\_rust}}
\renewcommand{\footrulewidth}{0.3pt}
\renewcommand{\footrule}{\hbox to\headwidth{\color{rulegray}\leaders\hrule height \footrulewidth\hfill}}
% Brand mark footer: CaseMirror wordmark left, page number right (both burgundy)
\fancyfoot[L]{\footnotesize\textbf{\color{headergray}CaseMirror}}
\fancyfoot[R]{\footnotesize\color{headergray}\thepage}

% -- Page 1 style: footer only, no header rule --
\fancypagestyle{firstpage}{%
  \fancyhf{}%
  \renewcommand{\headrulewidth}{0pt}%
  \renewcommand{\footrulewidth}{0.3pt}%
  \renewcommand{\footrule}{\hbox to\headwidth{\color{rulegray}\leaders\hrule height 0.3pt\hfill}}%
  \fancyfoot[L]{\footnotesize\textbf{\color{headergray}CaseMirror}}%
  \fancyfoot[R]{\footnotesize\color{headergray}\thepage}%
}

% -- Column rule --
\setlength{\columnseprule}{0.15pt}
\renewcommand{\columnseprulecolor}{\color{rulegray}}

% -- Title --
\title{\vspace{-1.8em}\LARGE\bfseries\color{titlebg} Codebase Review: beads\textbackslash{}\_rust\vspace{-0.3em}}
\author{CaseMirror Research \\ \small\itshape\href{https://casemirror.ai}{\textcolor{headergray}{casemirror.ai}}}
\date{{\small\color{headergray} \today}}

\begin{document}

\maketitle
\thispagestyle{firstpage}
\vspace{-0.5em}


\textbf{Date:} 2026-03-03

\textbf{Scope:} Full codebase (src/)

\textbf{Reviewer:} Claude Opus 4.6 (automated)


\subsection{Summary}

\begin{itemize}
\setlength{\itemsep}{0pt}
\setlength{\parskip}{0pt}
  \item 2 CRITICAL findings (bugs)
  \item 3 IMPORTANT findings (2 bugs, 1 performance)
\end{itemize}

Overall: well-structured codebase with parameterized queries throughout, \texttt{unsafe\_code = "forbid"}, path allowlists, atomic file writes. No SQL injection or path traversal vectors found. Issues are concentrated in \texttt{src/storage/sqlite.rs} where some code paths bypass the \texttt{mutate} transaction pattern.


\subsection{CRITICAL}

\subsubsection{BUG-1: Cycle detection misses waits-for dependencies}

\begin{itemize}
\setlength{\itemsep}{0pt}
\setlength{\parskip}{0pt}
  \item \textbf{File:} \texttt{src/storage/sqlite.rs:448}
  \item \textbf{Category:} Logic bug
  \item \textbf{Severity:} Critical
\end{itemize}

\texttt{check\_cycle} with \texttt{blocking\_only=true} filters to \texttt{('blocks', 'parent-child', 'conditional-blocks')} but omits \texttt{'waits-for'}, despite \texttt{DependencyType::is\_blocking()} and \texttt{rebuild\_blocked\_cache\_impl} both treating \texttt{waits-for} as blocking.


A \texttt{waits-for} cycle (A waits-for B, B waits-for A) will not be detected and will be inserted, corrupting the dependency graph. The blocked cache will then loop indefinitely or produce wrong results when the graph is traversed.


\textbf{Evidence:}


The model definition (\texttt{src/model/mod.rs} line 254-258):


\begin{lstlisting}
pub const fn is_blocking(&self) -> bool {
    matches!(
        self,
        Self::Blocks | Self::ParentChild | Self::ConditionalBlocks | Self::WaitsFor
    )
}
\end{lstlisting}

The blocked cache builder (\texttt{sqlite.rs} line 1494) also includes it:


\begin{lstlisting}
WHERE d.type IN ('blocks', 'conditional-blocks', 'waits-for')
\end{lstlisting}

But \texttt{check\_cycle} (line 448) omits it:


\begin{lstlisting}
let type_filter = if blocking_only {
    "AND type IN ('blocks', 'parent-child', 'conditional-blocks')"
} else {
    ""
};
\end{lstlisting}

\textbf{Fix:} Add \texttt{'waits-for'} to the type filter string on line 448:


\begin{lstlisting}
"AND type IN ('blocks', 'parent-child', 'conditional-blocks', 'waits-for')"
\end{lstlisting}

\subsubsection{BUG-2: add\textbackslash{}\_dependency cycle check is outside the transaction (TOCTOU)}

\begin{itemize}
\setlength{\itemsep}{0pt}
\setlength{\parskip}{0pt}
  \item \textbf{File:} \texttt{src/storage/sqlite.rs:1965-1975}
  \item \textbf{Category:} Race condition
  \item \textbf{Severity:} Critical
\end{itemize}

The cycle guard runs before \texttt{self.mutate(...)}, so it executes outside the \texttt{BEGIN IMMEDIATE} transaction. Two concurrent agents can simultaneously pass the check and both insert their dependency, creating a cycle.


\begin{lstlisting}
pub fn add_dependency(...) -> Result<bool> {
    // Cycle check runs HERE - no write lock held
    if let Ok(dt) = dep_type.parse::<DependencyType>()
        && dt.is_blocking()
        && self.would_create_cycle(issue_id, depends_on_id, true)?
    {
        return Err(BeadsError::DependencyCycle { ... });
    }

    // Actual write happens HERE under BEGIN IMMEDIATE
    self.mutate("add_dependency", actor, |conn, ctx| { ... })
}
\end{lstlisting}

In a concurrent multi-agent environment (the stated use case), two agents can simultaneously call \texttt{add\_dependency("A","B","blocks")} and \texttt{add\_dependency("B","A","blocks")}. Both read a cycle-free graph, both pass the check, and both insert their dependency, creating an A to B to A cycle.


\textbf{Compare:} \texttt{update\_issue}'s claim guard correctly re-checks the assignee inside the \texttt{mutate} closure (line 498-523) for this exact reason.


\textbf{Fix:} Move the cycle check inside the \texttt{mutate} closure so it executes under the exclusive write lock.


\subsection{IMPORTANT}

\subsubsection{BUG-3: set\textbackslash{}\_config/set\textbackslash{}\_metadata non-atomic DELETE+INSERT}

\begin{itemize}
\setlength{\itemsep}{0pt}
\setlength{\parskip}{0pt}
  \item \textbf{File:} \texttt{src/storage/sqlite.rs:2964-2976, 3358-3368}
  \item \textbf{Category:} Data integrity
  \item \textbf{Severity:} Important
\end{itemize}

Both perform DELETE then INSERT on the raw connection with no wrapping transaction. If INSERT fails (disk full, I/O error, process killed), the key is silently deleted with no recovery. Same pattern in \texttt{set\_export\_hash} and \texttt{set\_export\_hashes}.


\begin{lstlisting}
pub fn set_config(&mut self, key: &str, value: &str) -> Result<()> {
    // Step 1: DELETE (committed immediately in autocommit mode)
    self.conn.execute_with_params(
        "DELETE FROM config WHERE key = ?",
        &[SqliteValue::from(key)],
    )?;
    // Step 2: INSERT - if this fails, the key is permanently gone
    self.conn.execute_with_params(
        "INSERT INTO config (key, value) VALUES (?, ?)",
        &[SqliteValue::from(key), SqliteValue::from(value)],
    )?;
    Ok(())
}
\end{lstlisting}

Note: \texttt{set\_metadata\_in\_tx} (line 3520) exists as a transactional version and is used from import paths. But the public \texttt{set\_config} and \texttt{set\_metadata} methods are not transactional.


\textbf{Fix:} Wrap in an explicit transaction or route through \texttt{self.mutate(...)}.


\subsubsection{BUG-4: rebuild\textbackslash{}\_blocked\textbackslash{}\_cache uses plain BEGIN instead of BEGIN IMMEDIATE}

\begin{itemize}
\setlength{\itemsep}{0pt}
\setlength{\parskip}{0pt}
  \item \textbf{File:} \texttt{src/storage/sqlite.rs:1466}
  \item \textbf{Category:} Concurrency
  \item \textbf{Severity:} Important
\end{itemize}

All other write paths use \texttt{self.mutate(...)} with \texttt{BEGIN IMMEDIATE} + 5-retry exponential backoff. This path uses plain \texttt{BEGIN}, creating a deferred transaction that can fail with \texttt{SQLITE\_BUSY} when it tries to acquire the write lock for \texttt{DELETE FROM blocked\_issues\_cache}.


\begin{lstlisting}
pub fn rebuild_blocked_cache(&mut self, force_rebuild: bool) -> Result<usize> {
    if !force_rebuild {
        return Ok(0);
    }
    self.conn.execute("BEGIN")?;  // plain BEGIN, not BEGIN IMMEDIATE
    match Self::rebuild_blocked_cache_impl(&self.conn) { ... }
}
\end{lstlisting}

\textbf{Fix:} Use \texttt{self.mutate(...)} or \texttt{BEGIN IMMEDIATE} with retry logic.


\subsubsection{PERF-1: O(n\^{}2) visited-set check in collect\textbackslash{}\_descendant\textbackslash{}\_ids}

\begin{itemize}
\setlength{\itemsep}{0pt}
\setlength{\parskip}{0pt}
  \item \textbf{File:} \texttt{src/storage/sqlite.rs:2716}
  \item \textbf{Category:} Performance
  \item \textbf{Severity:} Important
\end{itemize}

BFS uses \texttt{Vec::contains()} as the visited check (O(n) per lookup). Called on every \texttt{--parent --recursive} query. \texttt{HashSet} is already imported at the top of the file.


\begin{lstlisting}
if !result.contains(&id) {  // O(n) linear scan
    queue.push_back(id.clone());
    result.push(id);
}
\end{lstlisting}

\textbf{Fix:} Replace with a \texttt{HashSet} for O(1) membership test:


\begin{lstlisting}
let mut visited = HashSet::new();
// ...
if visited.insert(id.clone()) {  // O(1) insert + membership test
    queue.push_back(id.clone());
    result.push(id);
}
\end{lstlisting}

\subsection{Positive Observations}

\begin{itemize}
\setlength{\itemsep}{0pt}
\setlength{\parskip}{0pt}
  \item Parameterized SQL queries everywhere -- no SQL injection vectors
  \item \texttt{#![forbid(unsafe\_code)]} enforced
  \item \texttt{mutate} pattern provides proper \texttt{BEGIN IMMEDIATE} + retry for most write paths
  \item Path allowlists prevent traversal attacks
  \item Atomic file writes for JSONL export
  \item The \texttt{check\_cycle} parameter ordering (confusing but correct) -- walks forward from \texttt{depends\_on\_id} to find \texttt{issue\_id}
\end{itemize}


\balance
\end{document}
